\definecolor{LightRed}{rgb}{1,0.5,0.5}
\begin{abstract}
In this work, we introduce long-video masked-embedding autoencoders (LV-MAE), a self-supervised learning framework for long video representation.
Our approach treats short- and long-span dependencies as two separate tasks.
Such decoupling allows for a more intuitive video processing where short-span spatiotemporal primitives are first encoded and are then used to capture long-range dependencies across consecutive video segments. 
To achieve this, we leverage advanced off-the-shelf multimodal encoders to extract representations from short segments within the long video, followed by pre-training a masked-embedding autoencoder capturing high-level interactions across segments.
LV-MAE is highly efficient to train and enables the processing of much longer videos by alleviating the constraint on the number of input frames.
Furthermore, unlike existing methods that typically pre-train on short-video datasets, our approach offers self-supervised pre-training using long video samples (\eg, 20+ minutes video clips) at scale.
Using LV-MAE representations, we achieve state-of-the-art results on three long-video benchmarks -- LVU, COIN, and Breakfast -- employing only a simple classification head for either attentive or linear probing.
Finally, to assess LV-MAE pre-training and visualize its reconstruction quality, we leverage the video-language aligned space of short video representations to monitor LV-MAE through video-text retrieval. Code is available at \href{https://github.com/amazon-science/lv-mae}{\textcolor{LightRed}{\texttt{https://github.com/amazon-science/lv-mae}}}.

\end{abstract}